\section{Obligation as Physical Scar}

For a Runic Warrior, obligation is not a number in a ledger. It is a scar on the hand, a brand on the chest, a missing finger, a limp that never quite heals. The covenant writes itself into flesh, and the flesh remembers what the mind might wish to forget.

This chapter provides mechanics and roleplay guidance for tracking Obligation as physical marks, as well as advice for Game Masters on using scars as narrative tools.

\subsection{The Visible Ledger}

A traditional Runekeeper's Obligation is invisible. Only the Patron knows the true balance. A Runic Warrior's Obligation is visible to anyone who looks.

\subsubsection{Mechanics}

When a Runic Warrior marks Obligation (from invoking a Rite, Pushing, or failing to meet a Patron's demand), the player should describe a physical change to their character. This change persists until the Obligation is cleared.

\begin{itemize}
\item \textbf{Minor Obligation (1-2 segments):} A small scar. A bruise that does not fade. A callus in a new place. A tooth that aches when the Patron's name is spoken.
\item \textbf{Moderate Obligation (3-4 segments):} A visible scar on the face or hands. A limp that worsens in cold weather. A tattoo that appears overnight. A brand that smokes when the Patron is angry.
\item \textbf{Significant Obligation (5-6 segments):} A missing finger. An eye that clouds. A patch of skin that turns to stone or bark. A second shadow that moves independently.
\item \textbf{Critical Obligation (7+ segments):} A limb that no longer feels pain. A face that reflections avoid. A heart that beats in a different rhythm. A name that others forget moments after hearing it.
\end{itemize}

These changes are not merely cosmetic. They affect how others perceive the Runekeeper.

\begin{itemize}
\item A scarred face may grant +1 die to Intimidation but -1 die to Persuasion.
\item A limp may reduce movement speed but grant advantage on checks to resist being knocked prone.
\item A missing finger may make fine manipulation harder but serve as a constant reminder of the debt.
\end{itemize}

\subsubsection{Clearing Obligation}

When a Runic Warrior clears Obligation (through service, sacrifice, or atonement), the physical mark does not simply vanish. It fades. A scar becomes a pale line. A brand becomes a smudge. A missing finger does not regrow—but the debt is no longer owed.

The character may choose to keep the mark as a reminder. Many do.

\subsection{Roleplaying Physical Obligation}

A scar is not a number. It is a story. Use these prompts to deepen roleplay.

\subsubsection{Questions for the Player}

\begin{itemize}
\item Which scar was the first? What Rite caused it? Did you know the cost before you invoked?
\item Has anyone asked about your marks? How did you answer? Did you lie, or tell the truth?
\item Do you hide your scars, or display them? Why?
\item Is there a scar you regret? A debt you wish you had never incurred?
\item Do the scars itch? Ache? Whisper? What do they say?
\end{itemize}

\subsubsection{Questions for the Game Master}

\begin{itemize}
\item Does an NPC recognize a particular scar as the mark of a specific Patron? How do they react?
\item Does a Chain-Lantern prefect use visible Obligation as evidence of unlawful covenant?
\item Does a Daughter of the Cord offer to help a Bloodsworn hide her scars—for a price?
\item Does a sentient Thiasos comment on a new scar? Does it approve? Disapprove?
\item Does a Patron ever touch a scar directly? What does that feel like?
\end{itemize}

\subsection{Sample Scar Tables}

Use these tables to generate physical marks of Obligation randomly, or for inspiration.

\subsubsection{Patron-Specific Scars}

\begin{itemize}
\item \textbf{Oath of Flame \& Light:} A burn that never heals. A palm that glows faintly in darkness. A tongue that tastes of ash when lying.
\item \textbf{Oath of Mercy \& Grace:} A scar over the heart. Hands that are always warm. A shadow that reaches toward the wounded.
\item \textbf{Mykkiel:} A brand shaped like a gavel. A finger that taps involuntarily when someone lies. A voice that echoes in empty rooms.
\item \textbf{Inquisitor Prime:} A censer mark on the forehead. Eyes that see corruption as a green haze. A cough that smells of brimstone.
\item \textbf{The Sealed Gate:} A scar in the shape of a door. Skin that is cold to the touch. A limp that worsens at thresholds.
\item \textbf{Grimmir:} Bark growing on the knuckles. Hair that turns green in spring. A scent of rain that precedes the Runekeeper.
\item \textbf{The Carrion-King:} A patch of skin that flakes like dead leaves. Breath that smells of compost. A single tooth that falls out and regrows each season.
\item \textbf{Clockwork Monad:} A gear pattern branded into the palm. An eye that ticks when focusing. Fingers that move in perfect sequence, never fumbling.
\item \textbf{Varnek Karn:} A scar that looks like a rune. A missing fingernail. A shadow that sometimes shows a different ancestor.
\item \textbf{The High Mother:} Tattoos that move when the Runekeeper sleeps. A hollowed bird's feather woven into the hair. A heartbeat that can be felt, not heard.
\item \textbf{Khemesh:} Skin that is always damp. Eyes that reflect light like deep water. A voice that sounds like stones grinding.
\item \textbf{Xhak'Thul:} Hair that stands on end before a storm. A scar that crackles when touched. A laugh that sounds like distant thunder.
\item \textbf{Ikasha:} A shadow that moves independently. Eyes that are hard to focus on. A whisper that follows the Runekeeper.
\item \textbf{The Silent Choir:} Lips that are sewn shut in a dream, then unsewn by morning. A scar across the throat. A voice that is never quite their own.
\item \textbf{Aveh:} A brand that looks like a broken chain. Skin that is always wind-chapped. A scent of ozone after a Rite.
\item \textbf{Inaea:} Threads visible beneath the skin. A scar that weaves across the body like a web. Fingers that are always slightly sticky.
\item \textbf{Isoka:} Patches of scale on the arms. A forked tongue that appears in dreams. A shed skin left on the pillow.
\item \textbf{Malachai:} A silver coin embedded in the flesh. Teeth that are sharp. A hunger that never quite fades.
\item \textbf{The Traveler:} A scar on the sole of the foot. A compass rose branded on the wrist. A shadow that always points north.
\item \textbf{The Witness:} A mirror-shard scar on the cheek. Eyes that reflect truth. A memory that cannot be forgotten.
\item \textbf{The Ninth:} A fractal scar that spirals. Eyes that see patterns in chaos. A voice that sometimes speaks in equations.
\item \textbf{Lunara:} A crescent brand on the forehead. Skin that pales under moonlight. A shadow that lengthens as the moon waxes.
\item \textbf{The Pale Shepherd:} A grey handprint on the shoulder. Breath that fogs in warm rooms. A bell that rings only for the dying.
\item \textbf{The Savage Heart:} A claw mark across the ribs. Teeth that are sharper than they should be. A growl that emerges unbidden.
\item \textbf{Lethai Ghe'hai:} A mask-shaped scar. Silk threads visible beneath the skin. A face that is hard to remember.
\end{itemize}

\subsection{Scars as Narrative Tools}

A scar is not a punishment. It is a prompt.

\subsubsection{Foreshadowing}

A scar that aches before a Rite backfires. A brand that warms when the Patron is near. A missing finger that itches when a debt is about to be called. Use physical marks to signal upcoming complications.

\subsubsection{Recognition}

An NPC recognizes a scar as the mark of a specific Patron—or a specific crime. A Chain-Lantern prefect has seen that brand before. A Daughter of the Cord knows that tattoo pattern. A mercenary captain lost a finger to a warrior with that same scar.

\subsubsection{Memory}

A scar holds memory. A Runic Warrior can touch a scar and recall the Rite that caused it, the battle where it was earned, the Patron who demanded payment. Use this for flashbacks or information gathering.

\subsubsection{Connection}

Two Runic Warriors may share similar scars—marks from the same Patron, the same Rite, the same battle. This creates an instant bond. It can also create suspicion.

\subsubsection{Sacrifice}

A Runic Warrior may offer a scar as payment. Not the healing of the wound—the scar itself. The memory of the pain. The story of the mark. The Patron accepts, and the scar fades. But the Runekeeper knows: something else is now written in its place.

\subsection{Optional Rule: Scarred Resilience}

Some Runic Warriors learn to draw strength from their scars. This optional talent represents that mastery.

\textbf{Scarred Resilience (4 XP)}

\textit{Requirement: Body 3+, at least three visible scars from Obligation.}

Once per session, when you would take Harm from a physical attack, you may instead convert that Harm into a new scar. The attack leaves a mark, but the pain is deferred. You do not take the Harm now.

However, the scar remains. It is visible. It is a debt recorded in flesh. The GM may call upon that scar later—as a complication, a clue, or a cost.

\subsection{Closing}

A scar is not a punishment. It is a receipt. It is a story written in flesh. It is a reminder that the covenant is real, that the debt was paid, that the warrior survived.

The thread held. The water carried. That is not a Rite. That is grace.
